\documentclass[a4paper]{article}

% Basic packages
\usepackage[fontset=fandol]{ctex}
\usepackage{amsmath, amssymb}
\usepackage{graphicx}
\usepackage[margin=2.5cm]{geometry}
\usepackage[most]{tcolorbox}
\usepackage{etoolbox}
\usepackage{listings}
\usepackage{hyperref}
\usepackage{booktabs}
\usepackage{subcaption}
\usepackage{float}
\usepackage{tikz}
\IfFileExists{pgfplots.sty}{
    \usepackage{pgfplots}
    \pgfplotsset{compat=1.18}
}{}

% ===== Highlight boxes =====
\newtcolorbox{knowledgebox}[1]{
    enhanced, colback=blue!5!white, colframe=blue!75!black, colbacktitle=blue!75!black,
    coltitle=white, fonttitle=\bfseries, title=#1, attach boxed title to top left={yshift=-2mm, xshift=2mm},
    boxrule=1pt, sharp corners
}

\newtcolorbox{importantbox}[1]{
    enhanced, colback=yellow!10!white, colframe=yellow!80!black, colbacktitle=yellow!80!black,
    coltitle=black, fonttitle=\bfseries, title=#1, sharp corners
}

\newtcolorbox{warningbox}[1]{
    enhanced, colback=red!5!white, colframe=red!75!black, colbacktitle=red!75!black,
    coltitle=white, fonttitle=\bfseries, title=#1, sharp corners
}

% ===== Code listings =====
\lstset{
    basicstyle=\ttfamily\small,
    keywordstyle=\color{blue},
    stringstyle=\color{red!60!black},
    commentstyle=\color{green!60!black},
    breaklines=true,
    frame=single,
    numbers=left,
    numberstyle=\tiny\color{gray},
    captionpos=b,
    extendedchars=false
}

% ===== Front-page metadata =====
\newcommand{\notetitle}{机器学习基本概念简介}
\newcommand{\notesubtitle}{李宏毅《机器学习》2025 · 第一节（上）}
\newcommand{\noteauthors}{基于李宏毅老师课程整理（lecture-to-notes 自动生成）}
\newcommand{\notedate}{\today}
\newcommand{\videochannel}{李宏毅深度学习课堂（Bilibili 搬运）}
\newcommand{\videopublishdate}{2025-08-06}
\newcommand{\videoduration}{约 50 分钟（49:59）}
\newcommand{\videourl}{https://www.bilibili.com/video/BV1TAtwzTE1S?p=1}
\newcommand{\videocoverpath}{cover.jpg}

\begin{document}

% ===== Title page =====
\begin{titlepage}
\centering
{\Large 课程笔记\par}
\vspace{1.2cm}
{\huge\bfseries \notetitle\par}
\ifdefempty{\notesubtitle}{}{
\vspace{0.3cm}
{\large \notesubtitle\par}
}
\vspace{0.5cm}
{\large \noteauthors\par}
\vspace{0.3cm}
{\large \notedate\par}
\vspace{1.0cm}

\ifdefempty{\videocoverpath}{
{\small 请在生成笔记时填入视频封面图的本地路径。\par}
}{
\includegraphics[width=0.82\textwidth,height=0.40\textheight,keepaspectratio]{\videocoverpath}\par
}

\vfill
\begin{tcolorbox}[width=0.9\textwidth, colback=black!2!white, colframe=black!60, sharp corners]
\textbf{视频作者/频道}：\videochannel\par
\textbf{发布日期}：\videopublishdate\par
\textbf{视频时长}：\videoduration\par
\textbf{视频链接}：
\ifdefempty{\videourl}{
未填写
}{
\href{\videourl}{\nolinkurl{\videourl}}
}
\end{tcolorbox}
\end{titlepage}

\tableofcontents
\newpage

%% --- 正文内容开始 --- %%
%% 正文将在 Whisper 字幕转写完成、并完成图文三方验证后填充。
%% 预定结构（基于 contact sheet 视觉分析，待字幕核对后定稿）：

\section{引言：机器学习在寻找什么}

机器学习（Machine Learning）常被科普文章渲染得“玄之又玄”——仿佛机器一旦学会学习，就有了人工智能，接下来就要统治人类。但若用一句话概括它的本质，其实相当朴素。

\begin{importantbox}{机器学习的本质}
机器学习就是\textbf{让机器具备“找一个函数”的能力}。这个函数往往极其复杂，复杂到人类无法用手写的方程式直接表达，于是我们寄希望于借助机器的力量，从资料中把它\textbf{自动}找出来。
\end{importantbox}

以语音辨识为例：我们想要的是一个函数 $f$，输入是一段声音讯号，输出是这段声音对应的文字内容。这样的函数无比复杂，绝不是人能手写出来的方程式。

\begin{figure}[H]
\centering
\includegraphics[width=0.72\textwidth]{figures/fig01_ml_function.png}
\caption{机器学习的核心是寻找一个人类难以手写的复杂函数。图中以语音辨识为例：函数输入声音波形，输出文字“How are you”\protect\footnotemark}
\end{figure}
\footnotetext{视频画面时间区间：00:02:00--00:02:15。}

这样的任务比比皆是，它们的共同点都是“找一个复杂的函数”：

\begin{itemize}
    \item \textbf{影像辨识}：$f(\text{一张图片}) = \text{图片里的内容}$；
    \item \textbf{下围棋（AlphaGo）}：$f(\text{棋盘上黑子与白子的位置}) = \text{下一步应落子的位置}$。
\end{itemize}

\begin{knowledgebox}{为什么要“学”出函数，而不是直接写}
像“声音讯号 $\to$ 文字”“棋盘 $\to$ 落子”这样的函数，输入输出关系极度复杂，人类没有能力把它显式地写成方程式。机器学习的价值，正是在人写不出函数时，凭借资料与算力把它逼近出来。
\end{knowledgebox}

\subsection{本章小结}
机器学习 $\approx$ 寻找函数。当任务的输入输出关系复杂到无法手写时，就交给机器从资料中自动寻找这个函数——这是贯穿全课程的第一性原理。

\section{函数的三种类型}

随着要找的函数不同，机器学习的任务可以分成几大类。这里先认识三个常见的专有名词。

\begin{itemize}
    \item \textbf{回归（Regression）}：函数的输出是一个\textbf{数值}（scalar）。例如预测明天中午的 PM2.5——输入今天的 PM2.5、平均温度、臭氧浓度等指标，输出一个数值。
    \item \textbf{分类（Classification）}：人类预先准备好若干\textbf{选项}（又称类别 class），函数从中选一个作为输出。例如 Gmail 侦测一封邮件“是 / 不是”垃圾邮件（两个选项）；AlphaGo 本质上也是分类——棋盘有 $19\times 19=361$ 个位置，相当于一道有 361 个选项的选择题。
    \item \textbf{结构化学习（Structured Learning）}：机器要产生一个\textbf{有结构}的物件，例如画一张图、写一篇文章。用拟人化的说法，就是让机器学会“创造”。
\end{itemize}

\begin{figure}[H]
\centering
\includegraphics[width=0.66\textwidth]{figures/fig02_structured_learning.png}
\caption{结构化学习（Structured Learning）：让机器创造出有结构的产物（图片、文档）。它位于回归与分类（图左下角）之外，是更具挑战性的一类任务\protect\footnotemark}
\end{figure}
\footnotetext{视频画面时间区间：00:08:15--00:08:30。}

\begin{knowledgebox}{回归与分类之外的“黑暗大陆”}
很多教科书讲到 Regression 和 Classification 就停下了，仿佛机器学习的世界只有“两大洲”。但在它们之外还有一块常被忽视的领域——Structured Learning，让机器创造有结构的输出。（李宏毅老师在此用《全职猎人》中的“黑暗大陆”作了一个比喻。）
\end{knowledgebox}

\subsection{本章小结}
按输出形态，机器学习任务分为三类：回归（输出数值）、分类（从选项中选择）、结构化学习（创造有结构的物件）。同一个问题（如下围棋）既可看作分类，也启发我们：任务类型决定了我们要找的是什么样的函数。

\section{机器学习的三个步骤}
% 贯穿案例：预测某频道每日观看人数

确立了“机器学习 $=$ 找函数”之后，真正的问题是：机器\textbf{怎么}找出这个函数？李宏毅老师以“预测自己的 YouTube 频道隔天的总观看人数”为贯穿案例，把找函数的过程拆解成三个步骤。

\subsection{第一步：带有未知参数的函数（Model）}

第一步，我们先\textbf{猜测}函数的数学形式——这个猜测来自我们对问题的理解，即领域知识（Domain Knowledge）。对于观看人数预测，一个最初步的猜测是：
$$ y = b + w\,x_1 $$
其中各符号的含义为：
\begin{itemize}
    \item $y$：要预测的目标，例如 2 月 26 日当天的总观看人数（今天尚未过完，所以未知）；
    \item $x_1$：已知的输入，即\textbf{特征}（feature）——前一天（2 月 25 日）的观看人数；
    \item $w$：与特征相乘的未知参数，称为\textbf{权重}（weight）；
    \item $b$：直接相加的未知参数，称为\textbf{偏置}（bias）。
\end{itemize}

\begin{figure}[H]
\centering
\includegraphics[width=0.82\textwidth]{figures/fig03_step1_model.png}
\caption{第一步——写出带有未知参数的函数 $y=b+wx_1$。其中 $x_1$ 是特征（前一天观看数），$w$、$b$ 是待从资料中学出的未知参数（权重与偏置）\protect\footnotemark}
\end{figure}
\footnotetext{视频画面时间区间：00:14:45--00:15:00。}

这种带有未知参数（unknown parameters）的函数，在机器学习中就叫做\textbf{模型}（model）。这里之所以猜测 $y=b+wx_1$，背后的直觉是：今天的观看数大概会和昨天的有关联，但不会一模一样，于是把昨天的数值乘上 $w$ 再加一个修正项 $b$。

\begin{knowledgebox}{Domain Knowledge 用在哪里}
经常有人说“做机器学习需要一些领域知识”。它最主要的用武之地，正是在第一步\textbf{写出带有未知参数的函数}时——你对问题了解越深，越能猜出一个合理的函数形式（model）。
\end{knowledgebox}

\subsection{第二步：从训练资料定义损失（Loss）}

第二步，定义一个叫做\textbf{损失}（Loss）的函数 $L(b, w)$。注意它本身也是一个函数，但它的\textbf{输入是模型的参数} $b$、$w$，输出是一个数值，用来衡量“这一组参数好不好”。

\begin{figure}[H]
\centering
\includegraphics[width=0.82\textwidth]{figures/fig04_loss_example.png}
\caption{损失的计算示例：固定 $b=0.5\mathrm{k}$、$w=1$ 时，用 1 月 1 日的真实值 $4.8\mathrm{k}$ 预测 1 月 2 日得 $5.3\mathrm{k}$，而真实值（标签）$\hat{y}=4.9\mathrm{k}$\protect\footnotemark}
\end{figure}
\footnotetext{视频画面时间区间：00:18:30--00:18:45。}

损失要用\textbf{训练资料}来计算。本例的训练资料是该频道 2017/01/01 至 2020/12/31 共三年每天的真实观看数。具体做法：先固定一组参数（如 $b=0.5\mathrm{k}, w=1$），模型变为 $y=0.5\mathrm{k}+x_1$；把 1 月 1 日的真实值 $4.8\mathrm{k}$ 代入，预测 1 月 2 日为 $5.3\mathrm{k}$，而 1 月 2 日的真实值（标签 label）是 $4.9\mathrm{k}$，二者的差距为
$$ e_1 = \lvert y - \hat{y}\rvert = \lvert 4.9\mathrm{k} - 5.3\mathrm{k}\rvert = 0.4\mathrm{k}. $$

对三年里的每一天都这样算出误差 $e_n$，再取平均，就得到损失：
$$ L = \frac{1}{N}\sum_{n} e_n $$
\begin{itemize}
    \item $N$：训练资料的笔数（三年约 $365\times 3$ 笔）；
    \item $e_n$：第 $n$ 笔资料上预测值与真实值的差距；
    \item $L$ 越大表示这组参数越差，$L$ 越小表示越好。
\end{itemize}

\begin{figure}[H]
\centering
\includegraphics[width=0.82\textwidth]{figures/fig05_loss_formula.png}
\caption{损失的一般式 $L=\frac{1}{N}\sum_n e_n$，以及单笔误差 $e$ 的两种常见定义：绝对值（MAE）与平方（MSE）\protect\footnotemark}
\end{figure}
\footnotetext{视频画面时间区间：00:21:45--00:22:00。}

单笔误差 $e$ 的计算方式不止一种：
\begin{itemize}
    \item 取绝对值 $e = \lvert y - \hat{y}\rvert$，对应的损失称为\textbf{平均绝对误差}（MAE, Mean Absolute Error）；
    \item 取平方 $e = (y - \hat{y})^2$，对应\textbf{均方误差}（MSE, Mean Square Error）。
\end{itemize}

\begin{knowledgebox}{MAE、MSE 与交叉熵怎么选}
MAE 与 MSE 之间存在微妙差别，选用哪一种取决于任务需求和你对问题的理解，并非随意。特别地，若 $y$ 与 $\hat{y}$ 都是几率分布，则常改用\textbf{交叉熵}（cross-entropy）来衡量差距。
\end{knowledgebox}

如果把每一组参数 $(w, b)$ 都试一遍、各自算出损失 $L$，就能把结果画成一张\textbf{误差曲面}（error surface）等高线图。

\begin{figure}[H]
\centering
\includegraphics[width=0.72\textwidth]{figures/fig06_error_surface.png}
\caption{误差曲面（error surface）：横轴 $w$、纵轴 $b$，颜色表示损失大小。越偏红色 $L$ 越大（参数越差），越偏蓝色 $L$ 越小（参数越好）\protect\footnotemark}
\end{figure}
\footnotetext{视频画面时间区间：00:24:30--00:24:45。}

在这张真实数据画出的等高线图上，损失最小（最蓝）的位置大约在 $w\approx 1$、$b$ 取一个较小的值附近。这其实符合直觉：用前一天的观看数预测隔天，二者往往差不多，因此当 $w\approx 1$、$b$ 仅作小幅修正时，预测最准。

\begin{importantbox}{误差曲面 error surface}
把每一组参数对应的损失可视化为等高线，便得到 error surface。它直观地告诉我们“往哪个方向走、参数怎么调”才能让损失变小——而这正是第三步“最佳化”要解决的问题。
\end{importantbox}
\subsection{第三步：最佳化（Optimization）}

第三步要解一个\textbf{最佳化}（optimization）问题：找出让损失最小的那组参数，记作 $w^*, b^*$：
$$ w^*,\, b^* = \arg\min_{w,\,b} L. $$
本课程使用的最佳化方法只有一种——\textbf{梯度下降}（Gradient Descent）。

为便于理解，先假设只有一个参数 $w$，此时 error surface 退化成一条一维曲线。梯度下降的做法是：

\begin{enumerate}
    \item 随机选取一个初始值 $w^0$；
    \item 计算损失对 $w$ 的微分（即该点切线的斜率）$\left.\dfrac{\partial L}{\partial w}\right|_{w=w^0}$；
    \item 沿着“让损失下降”的方向，把 $w$ 挪动一步。
\end{enumerate}

\begin{figure}[H]
\centering
\includegraphics[width=0.82\textwidth]{figures/fig07_gradient_descent.png}
\caption{梯度下降：从随机初始点 $w^0$ 出发，按 $w^{1}\leftarrow w^{0}-\eta\frac{\partial L}{\partial w}$ 反复更新；右侧标出了全局最小（global minima）与可能困住它的局部最小（local minima）\protect\footnotemark}
\end{figure}
\footnotetext{视频画面时间区间：00:35:00--00:35:15。}

斜率为负（左高右低）就增大 $w$，斜率为正（左低右高）就减小 $w$。形象地说，就像一个人站在山坡上左右环视，朝更低的方向跨出一步。这一步跨多大，取决于两件事：当前斜率的大小，以及一个你自己设定的\textbf{学习率}（learning rate）$\eta$。更新公式为：
$$ w^{1} \leftarrow w^{0} - \eta\,\frac{\partial L}{\partial w}\bigg|_{w=w^{0}} $$
\begin{itemize}
    \item $\eta$：学习率，由人自行设定，控制每一步的幅度；
    \item $\partial L/\partial w$：当前位置的斜率（梯度）；
    \item 负号：保证朝损失\textbf{下降}的方向移动。
\end{itemize}

\begin{knowledgebox}{超参数 hyperparameter}
像学习率 $\eta$、以及“最多更新几次”这类需要\textbf{人为设定}、而非机器从资料中学出来的量，统称为\textbf{超参数}（hyperparameter）。它与“参数”（$w$、$b$，由机器学出）是两个不同层级的概念。
\end{knowledgebox}

反复迭代 $w^0\to w^1\to w^2\to\cdots$，通常在两种情况下停止：（1）达到预设的更新次数上限（“失去耐心”，次数本身也是超参数）；（2）微分恰好为 0，参数不再移动。

\begin{warningbox}{Local minima 是个“假议题”}
梯度下降一旦走到斜率为 0 处就会停下，因此可能停在并非全局最低的\textbf{局部最小}（local minima），而非\textbf{全局最小}（global minima）。教科书常把这说成梯度下降的致命缺陷。但李宏毅老师特别提醒：在真实的深度学习里，local minima 其实是个“假议题”——梯度下降真正棘手的难题另有其物，留待后续课程揭晓（伏笔见第四节：局部最小与鞍点）。
\end{warningbox}

推广到两个参数 $w$、$b$ 时，做法完全一样：分别计算 $\partial L/\partial w$ 与 $\partial L/\partial b$，再同步更新二者。这些微分在 PyTorch 等深度学习框架里由一行程式自动算出，无需手算。

\begin{figure}[H]
\centering
\includegraphics[width=0.78\textwidth]{figures/fig08_gd_two_params.png}
\caption{双参数梯度下降：在 $(w,b)$ 的误差曲面上，每一步沿 $(-\eta\frac{\partial L}{\partial w},\,-\eta\frac{\partial L}{\partial b})$ 方向移动，逐步逼近损失更小的区域\protect\footnotemark}
\end{figure}
\footnotetext{视频画面时间区间：00:39:30--00:39:45。}

在真实数据上跑梯度下降，得到最佳参数 $w^*=0.97$、$b^*\approx 0.1\mathrm{k}$，训练资料（2017–2020）上的损失 $L=0.48\mathrm{k}$。这与最初的猜测 $w\approx 1$ 相当接近——用前一天预测隔天，本就该有 $w\approx 1$。

\subsection{本章小结}

机器学习“找函数”的三个步骤至此完整：

\begin{importantbox}{机器学习的三个步骤（= 训练 Training）}
\begin{enumerate}
    \item \textbf{Model}：写出一个带未知参数的函数 $y=b+wx_1$；
    \item \textbf{Loss}：定义损失 $L(b,w)$，用训练资料衡量参数好坏；
    \item \textbf{Optimization}：用梯度下降找出 $w^*,b^*$ 使损失最小。
\end{enumerate}
三步合起来称为\textbf{训练}（training）。
\end{importantbox}

\begin{figure}[H]
\centering
\includegraphics[width=0.85\textwidth]{figures/fig09_ml_simple_summary.png}
\caption{“Machine Learning is so simple” —— 三步骤总览，以及本例求得的最佳参数 $w^*=0.97$、$b^*=0.1\mathrm{k}$、训练损失 $L=0.48\mathrm{k}$\protect\footnotemark}
\end{figure}
\footnotetext{视频画面时间区间：00:40:30--00:40:45。}

\section{从训练到预测：泛化的考验}

然而 $L=0.48\mathrm{k}$ 是在\textbf{已经知道答案}的训练资料上算出来的——用李宏毅老师的话说，这只是“自嗨”。我们真正关心的，是模型对\textbf{未来未知数据}的预测能力。于是用 2020 年底找到的函数，逐日预测 2021 年（1 月 1 日至 2 月 14 日）的观看数：在这些\textbf{没看过}的数据上，误差 $L'=0.58\mathrm{k}$，明显大于训练误差 $0.48\mathrm{k}$。

\begin{figure}[H]
\centering
\includegraphics[width=0.82\textwidth]{figures/fig10_train_vs_predict.png}
\caption{2021 年预测结果：横轴为时间，红线是真实观看人次、蓝线是模型预测。蓝线几乎就是红线“向右平移一天”\protect\footnotemark}
\end{figure}
\footnotetext{视频画面时间区间：00:43:15--00:43:30。}

\begin{warningbox}{训练误差 $\neq$ 测试误差}
在见过的资料上表现好，不代表在没见过的资料上也好。我们真正追求的是后者，即\textbf{泛化能力}。这是机器学习中必须时刻警惕的区别。
\end{warningbox}

仔细观察预测曲线还能发现：蓝线几乎是红线右移一天——因为这个模型本质上就是“拿前一天预测隔天”（$w\approx 1$）。而真实数据有一个明显的\textbf{7 天周期}：每周五、周六观看数特别低（周末大家都出去玩，没人学机器学习）。

\section{让模型更有弹性：线性模型}

既然数据有 7 天周期，只看前一天的模型显然太弱。于是把前 7 天都纳入考虑——这次对模型的改良，又一次来自领域知识（domain knowledge）：
$$ y = b + \sum_{j=1}^{7} w_j\,x_j $$
其中 $x_j$ 是 $j$ 天前的观看数，$w_j$ 是它对应的权重。

\begin{figure}[H]
\centering
\includegraphics[width=0.78\textwidth]{figures/fig11_multi_feature.png}
\caption{从单特征 $y=b+wx_1$ 扩展到多特征 $y=b+\sum_{j=1}^{7}w_j x_j$：把前若干天的观看数都作为特征纳入考虑\protect\footnotemark}
\end{figure}
\footnotetext{视频画面时间区间：00:47:00--00:47:15。}

继续增加纳入考虑的天数，训练与测试损失的变化如下表：

\begin{table}[H]
\centering
\begin{tabular}{ccc}
\toprule
纳入的天数 & 训练损失 $L$（2017--2020） & 测试损失 $L'$（2021） \\
\midrule
1 天   & 0.48k & 0.58k \\
7 天   & 0.38k & 0.49k \\
28 天  & 0.33k & 0.46k \\
56 天  & 0.32k & 0.46k \\
\bottomrule
\end{tabular}
\caption{考虑不同天数的线性模型，在训练与测试资料上的损失对比。}
\end{table}

可以看到：纳入越多天，训练损失持续下降；但\textbf{测试损失}在 28 天之后就停在约 $0.46\mathrm{k}$ 不再改善——靠“多看几天”来提升，已经到了极限。

\begin{importantbox}{线性模型 Linear Model}
上面这些模型都是“特征乘以权重、再加偏置”的形式：
$$ y = b + \sum_{j} w_j\,x_j. $$
它们有一个共同的名字——\textbf{线性模型}（linear model）。
\end{importantbox}

\subsection{本章小结}
借助 domain knowledge 引入更多特征（前 $N$ 天），线性模型把测试误差从 $0.58\mathrm{k}$ 降到约 $0.46\mathrm{k}$，但很快触及瓶颈。这暗示：线性模型的表达能力存在天然上限。

\section{总结与延伸}

\textbf{讲者的收束}：本节用一个贯穿始终的例子（预测频道观看数），完整走过了机器学习的三个步骤，并亲眼看到线性模型的能力极限。李宏毅老师在结尾点明下一步——“看看怎么把 linear model 做得更好”。这正是通往更复杂模型、乃至深度学习的起点。

\textbf{结构化提炼}：

\begin{itemize}
    \item \textbf{一条主线}：机器学习 $=$ 找一个函数；而“找函数”被拆解为三步骤——Model $\to$ Loss $\to$ Optimization。
    \item \textbf{两个关键区分}：参数（$w,b$，机器学出）与超参数（$\eta$、迭代次数等，人为设定）；训练误差与测试误差（真正重要的是后者，即泛化）。
    \item \textbf{一个伏笔}：梯度下降的 local minima 被老师称为“假议题”，真正的难题留待后续揭晓。
    \item \textbf{一处局限}：线性模型即便不断增加特征，也存在天花板，需要更有弹性的模型。
\end{itemize}

\textbf{下一步}：第一节（下）将引入 Sigmoid / ReLU，把线性模型扩展成能逼近任意曲线的函数，进而引出“类神经网络”与深度学习的基本概念。

%% --- 正文内容结束 --- %%

\end{document}
